% answers_short.tex —— 2013 年·天津卷·理科 · 简短答案速查
% 由 tools/make_short_answers.py 自 answers.tex 的 \gksolhead/\gkanswer 提取，
% 并与 2013/2013_数学真题_天津理_答案要点.md 逐题比对后生成，勿手改（改源 MD 后重跑）。

\begin{gknotice}
\textbf{使用说明}：本节只给最终答案便于快速核对，完整推导见解析版。答案为编写者独立推导，\textbf{非官方标准答案}。
\end{gknotice}

\gkpart{选择题}
\begin{center}\begin{tabular}{ccccccccc}
\toprule
题号 & 1 & 2 & 3 & 4 & 5 & 6 & 7 & 8 \\
\midrule
答案 & D & A & B & C & C & C & B & A \\
\bottomrule
\end{tabular}\end{center}

\gkpart{填空题}
\begin{itemize}[leftmargin=2.6em]
  \item[\textbf{9.}] $1+2\mathrm i$。展开得 $(a-1)+(a+1)\mathrm i=b\mathrm i$，故 $a=1$，$b=2$。
  \item[\textbf{10.}] $15$。通项为 $\mathrm C_6^k(-1)^kx^{6-\frac{3k}{2}}$，令 $6-\dfrac{3k}{2}=0$ 得 $k=4$，常数项为 $\mathrm C_6^4=15$。
  \item[\textbf{11.}] $2\sqrt3$。圆的直角坐标方程为 $(x-2)^2+y^2=4$，圆心 $C(2,0)$；$P$ 的直角坐标为 $(2,2\sqrt3)$，故 $|CP|=2\sqrt3$。
  \item[\textbf{12.}] $\dfrac12$。设 $\overrightarrow{AB}=\boldsymbol a$，$\overrightarrow{AD}=\boldsymbol b$，由 $\overrightarrow{AC}\cdot\overrightarrow{BE}=-\dfrac{|\boldsymbol a|^2}{2}+\dfrac{|\boldsymbol a|}{4}+1=1$ 得 $|\boldsymbol a|=\dfrac12$。
  \item[\textbf{13.}] $\dfrac83$。切割线定理得 $EB=4$，$ED=9$；由弦切角与 $BD\parallel AC$ 得 $\triangle EAB\sim\triangle BCA$，故 $AB=4$，$BC=6$；再由 $\triangle FBD\sim\triangle FAC$ 得 $\dfrac{BF}{AF}=\dfrac{FD}{FC}=\dfrac54$，结合 $BF+FC=6$、$AD=6$ 解得 $FC=\dfrac83$。
  \item[\textbf{14.}] $-2$。$a>0$ 时最小值 $\dfrac54$（$a=\dfrac23$）；$a<0$ 时最小值 $\dfrac34$（$a=-2$）。故取 $a=-2$。
\end{itemize}

\gkpart{解答题}
\begin{itemize}[leftmargin=2.6em]
  \item[\textbf{15.}] (1) $\pi$。化简得 $f(x)=2\sqrt2\sin\left(2x-\dfrac\pi4\right)$，故最小正周期为 $\pi$；(2) 最大值 $2\sqrt2$，最小值 $-2$。在 $\left[0,\dfrac\pi2\right]$ 上 $2x-\dfrac\pi4\in\left[-\dfrac\pi4,\dfrac{3\pi}{4}\right]$，故 $f(x)\in[-2,2\sqrt2]$
  \item[\textbf{16.}] (1) $\dfrac67$。总取法 $\mathrm C_7^4=35$，不含编号 $3$ 的取法为 $\mathrm C_5^4=5$，故所求为 $1-\dfrac5{35}=\dfrac67$；(2) $P(X=1)=\dfrac1{35}$，$P(X=2)=\dfrac4{35}$，$P(X=3)=\dfrac{10}{35}=\dfrac27$，$P(X=4)=\dfrac{20}{35}=\dfrac47$；$E(X)=\dfrac{17}{5}$。计数分别为 $1,4,10,20$，总和 $35$，期望 $\dfrac{1+8+30+80}{35}=\dfrac{17}{5}$
  \item[\textbf{17.}] (1) 证明见解析。取 $A$ 为原点建系，$\overrightarrow{B_1C_1}=(-1,1,0)$，$\overrightarrow{CE}=(-1,-1,1)$，内积为 $0$；(2) $\dfrac{\sqrt{21}}7$。两平面法向量可取 $\boldsymbol n_1=(1,-3,-2)$，$\boldsymbol n_2=(2,-2,0)$，故 $\cos\theta=\dfrac2{\sqrt7}$，$\sin\theta=\dfrac{\sqrt{21}}7$；(3) $\sqrt2$。设 $M(t,t,1+t)$，由 $\dfrac{t}{\sqrt{3t^2+2t+1}}=\dfrac{\sqrt2}{6}$ 得 $15t^2-2t-1=0$，取 $t=\dfrac13$，故 $AM=\sqrt2$
  \item[\textbf{18.}] (1) $\dfrac{x^2}{3}+\dfrac{y^2}{2}=1$。由 $e=\dfrac{\sqrt3}{3}$ 得 $b^2=\dfrac23a^2$；通径长 $\dfrac{4\sqrt3}{3}$ 给出 $b^2=2$，$a^2=3$；(2) $k=\pm\sqrt2$。联立直线 $y=k(x+1)$ 与椭圆，韦达定理代入 $\overrightarrow{AC}\cdot\overrightarrow{DB}+\overrightarrow{AD}\cdot\overrightarrow{CB}=6-2(x_1x_2+y_1y_2)=8$，得 $k^2=2$
  \item[\textbf{19.}] (1) $a_n=\dfrac{3(-1)^{n-1}}{2^n}$（$n\in\mathbf N^*$）。由等差条件得 $q^2(4q^2-1)=0$，即 $q=\pm\dfrac12$；$q=\dfrac12$ 时为递减数列，舍去，故 $q=-\dfrac12$；(2) 最大项 $T_1=\dfrac56$，最小项 $T_2=-\dfrac7{12}$。$S_n=1-\left(-\dfrac12\right)^n$；偶数项 $T_n$ 递增，奇数项 $T_n$ 递减，且分别以 $0$ 为极限
  \item[\textbf{20.}] (1) 单调递减区间 $\left(0,e^{-\frac12}\right)$，单调递增区间 $\left(e^{-\frac12},+\infty\right)$。由 $f'(x)=x(2\ln x+1)$ 判号即得；(2) 证明见解析。$f$ 在 $\left(e^{-\frac12},+\infty\right)$ 上严格递增，值域为 $\left[-\dfrac1{2e},+\infty\right)$，而 $t>0>-\dfrac1{2e}$，故存在唯一解；另一段 $f<0$ 不产生解；(3) 证明见解析。由 $t>e^2=f(e)$ 得 $s>e$；比较 $f(\sqrt t)=\dfrac t2\ln t>t$ 得 $\sqrt t>s$，给出上界；比较 $f\left(t^{\frac25}\right)=\dfrac25t^{\frac45}\ln t<t$ 得 $t^{\frac25}<s$，给出下界。数值复核 $t=e^4,e^5,e^{10}$ 时比值分别约为 $0.4317$，$0.4247$，$0.4274$，均落在 $\left(\dfrac25,\dfrac12\right)$ 内
\end{itemize}
